% Source: https://artofproblemsolving.com/community/c3h3158474_2016_lmt_spring_guts_round_p13p24__lexington_mathematical_tournament
% Stable identity: aops-topic-3158474-q22
Repeated use of $2\cos^2\theta-1=\cos2\theta$ and the positive half-angle identity shows that a nest with $r$ roots, divided by $2$, equals $\cos\dfrac{\pi}{2^{r+1}}$. Thus $m=\cos\dfrac{\pi}{2^{2015}}$. Each application of $f_1$ doubles the angle, so $f_{2015}(m)=\cos\pi=\boxed{-1}$.
\begin{asy}
usepackage("amsmath");
size(12cm);
pair O=(0,0),A=dir(0),B=dir(30),C=dir(60);draw(circle(O,1),gray);draw(O--A^^O--B^^O--C);draw(B--(B.x,0)^^C--(C.x,0),dashed);draw(arc(O,.32,0,30),blue);draw(arc(O,.55,0,60),red);label("$O$",O,SW);label("$1$",O--C,NW);label("$\theta$",.4dir(15),SE);label("$2\theta$",.62dir(45),NW);label("$\cos\theta$",(B.x/2,-.09),S);label("$\cos2\theta$",(C.x/2,-.35),S);label("Generic half-angle construction",(0,-1.3));label("$\cos\theta=\dfrac{\sqrt{2+2\cos2\theta}}{2}$",(0,-1.65));
\end{asy}
